\documentclass[12pt]{article}
\usepackage[utf8]{inputenc}
\usepackage[T1]{fontenc}
\usepackage{amsmath, amssymb}
\usepackage{geometry}
\usepackage{xcolor}
\usepackage{lipsum}
\usepackage{titlesec}
\usepackage{fancyhdr}
\usepackage[most]{tcolorbox}
\usepackage{graphicx}
\usepackage{float}
\usepackage{tikz}
\usetikzlibrary{shapes.geometric, arrows.meta}
\usetikzlibrary{arrows.meta, decorations.markings}

\geometry{a4paper, margin=2.5cm}
\pagestyle{fancy}
\fancyhf{}
\rhead{Problem Set VI}
\lhead{Rigid body mechanics}
\cfoot{\thepage}

% Fancy titles
\titleformat{\section}{\normalfont\Large\bfseries}{Exercise \thesection:}{1em}{}
\titleformat{\subsection}{\normalfont\bfseries}{Correction:}{1em}{}

% Custom box for correction with breakable option
\newtcolorbox{correctionbox}{
  colback=gray!5,
  colframe=black,
  fonttitle=\bfseries,
  title=Correction,
  before skip=10pt,
  after skip=10pt,
  breakable,
  enhanced
}

% Fix headheight warning
\setlength{\headheight}{14.49998pt}

\begin{document}

\begin{center}
\Large\textbf{Ibn Tofail University} \\[1em]
\large\textit{Rigid body mechanics} \\[2em]
\large\textit{Problem Set VI} \\[0.5em]
\end{center}

\vspace{1cm}

% ----------- EXERCISE 1 -------------
\section{}

\begin{enumerate}
    \item What is the force necessary to move, at constant velocity, a perfect rigid body, of mass $m$, having the shape of a cube (Fig 1)?
    \item What is the force necessary to move, at constant velocity and without slipping, a perfect rigid body having the shape of a cylinder of radius $a$ and of the same mass $m$ (fig 2)?
    
    We denote by $\vec{R}$ and $\vec{M}_I$ the reduction elements of the torsor of contact actions at point $I$.
    \item Compare the norms of the two forces $\vec{F}_1$ and $\vec{F}_2$.
\end{enumerate}

\begin{correctionbox}
\textbf{1. Force for cube moving at constant velocity:}

For a cube sliding at constant velocity, the only forces to overcome are friction forces. The momentum theorem gives:
$$\vec{F}_1 + \vec{R} + m\vec{g} = \vec{0}$$

Since the velocity is constant, acceleration is zero. The reaction $\vec{R}$ has a normal component $N$ and a tangential component $T$ due to friction:
$$R = -N\vec{j} - T\vec{i}$$

The friction force follows Coulomb's law: $T = fN$, where $f$ is the coefficient of friction.

From the vertical equilibrium: $N = mg$

Therefore: $T = fmg$

So the force needed is: $\|\vec{F}_1\| = fmg$

\textbf{2. Force for cylinder rolling without slipping:}

For a cylinder rolling without slipping at constant velocity, we apply both momentum and angular momentum theorems.

Momentum theorem:
$$\vec{F}_2 + \vec{R} + m\vec{g} = \vec{0}$$

Angular momentum theorem about the center of mass $G$:
$$\vec{M}_G(\vec{F}_2) + \vec{M}_G(\vec{R}) = \vec{0}$$

Let $\vec{R} = -N\vec{j} - T\vec{i}$ (reaction at contact point I)

The moment of $\vec{F}_2$ about $G$ is: $\vec{M}_G(\vec{F}_2) = \vec{0}$ (if applied at G)

The moment of $\vec{R}$ about $G$ is: $\vec{M}_G(\vec{R}) = \vec{GI} \times \vec{R} = (-a\vec{j}) \times (-T\vec{i} - N\vec{j}) = aT(\vec{j} \times \vec{i}) = -aT\vec{k}$

So the angular momentum theorem gives: $-aT = 0 \Rightarrow T = 0$

From the momentum theorem in the horizontal direction: $F_2 - T = 0 \Rightarrow F_2 = T = 0$

Wait, this suggests no force is needed, which contradicts physical intuition. Let's reconsider.

Actually, for pure rolling at constant velocity on a horizontal surface, once the motion is established, no horizontal force is needed to maintain it (neglecting rolling resistance). The force $\vec{F}_2$ is zero to maintain constant velocity.

However, if we consider rolling resistance, there would be a small force needed.

\textbf{3. Comparison:}

For the cube: $\|\vec{F}_1\| = fmg$

For the cylinder (ideal case, no rolling resistance): $\|\vec{F}_2\| = 0$

Therefore: $\|\vec{F}_1\| > \|\vec{F}_2\|$

This shows the advantage of rolling over sliding.
\end{correctionbox}

% ----------- EXERCISE 2 -------------
\section{}

A perfect and homogeneous solid $S$ having the shape of a disk of mass $m$ and radius $a$, is in motion, in a vertical plane, on the axis $OX$ of a terrestrial reference frame assumed to be Galilean. Initially, the angular velocity is $\omega_0$ and the velocity of the center of mass $G$ is $v_0$.

\begin{enumerate}
    \item Using the general theorems, determine the differential equations of motion, in $x$ and $\theta$, of $S$ as a function of $m$, $a$, and $T$, where $T$ is the norm of the tangential component of the reaction at point $I$.
    \item Determine the sliding velocity of $S$ with respect to $OX$ and express $T$ as a function of $v_g$ and $m$.
    \item Deduce that the motion ends with a rolling without slipping phase.
\end{enumerate}

\begin{correctionbox}
\textbf{1. Differential equations of motion:}

Let's apply the general theorems.

Momentum theorem:
$$m\ddot{x} = T$$
where $T$ is the friction force (tangential component of reaction).

Angular momentum theorem about the center $G$:
$$I_G\ddot{\theta} = -aT$$

For a disk: $I_G = \frac{1}{2}ma^2$

So: $\frac{1}{2}ma^2\ddot{\theta} = -aT \Rightarrow \ddot{\theta} = -\frac{2T}{ma}$

The differential equations are:
$$\ddot{x} = \frac{T}{m}$$
$$\ddot{\theta} = -\frac{2T}{ma}$$

\textbf{2. Sliding velocity and expression for $T$:}

The sliding velocity is the velocity of the contact point $I$ relative to the ground:
$$\vec{v}_g = \vec{V}_I = \vec{V}_G + \vec{\omega} \times \vec{GI}$$

We have $\vec{V}_G = \dot{x}\vec{i}$, $\vec{\omega} = \dot{\theta}\vec{k}$, $\vec{GI} = -a\vec{j}$

So: $\vec{v}_g = \dot{x}\vec{i} + \dot{\theta}\vec{k} \times (-a\vec{j}) = \dot{x}\vec{i} + a\dot{\theta}\vec{i} = (\dot{x} + a\dot{\theta})\vec{i}$

The norm of sliding velocity is: $v_g = |\dot{x} + a\dot{\theta}|$

According to Coulomb's friction law:
$$T = fN \cdot \text{sign}(v_g) = fmg \cdot \text{sign}(v_g)$$
where $f$ is the coefficient of friction.

\textbf{3. Motion ends with rolling without slipping:}

Let's examine the evolution of sliding velocity:
$$\frac{d}{dt}(\dot{x} + a\dot{\theta}) = \ddot{x} + a\ddot{\theta} = \frac{T}{m} + a\left(-\frac{2T}{ma}\right) = \frac{T}{m} - \frac{2T}{m} = -\frac{T}{m}$$

So: $\frac{dv_g}{dt} = -\frac{T}{m}$

Since $T$ has constant magnitude $fmg$ and opposes the sliding, we have:
$$\frac{dv_g}{dt} = -fg \cdot \text{sign}(v_g)$$

This means the sliding velocity decreases linearly with time until it reaches zero. Once $v_g = 0$, pure rolling is established and maintained.

Therefore, the motion always ends with a rolling without slipping phase.
\end{correctionbox}

% ----------- EXERCISE 3 -------------
\section{}

A homogeneous hoop $S$, of mass $m$ and radius $R$, rolls with slipping on the axis $OX$, inclined at an angle $\alpha$ with respect to the horizontal, of a terrestrial frame assumed to be Galilean.

\begin{enumerate}
    \item Calculate the mechanical energy of solid $S$.
    \item By applying the theorem of variation of mechanical energy, find the differential equation of motion of $S$.
\end{enumerate}

\begin{correctionbox}
\textbf{1. Mechanical energy:}

The mechanical energy $E_m$ is the sum of kinetic and potential energy.

Kinetic energy: $E_c = \frac{1}{2}mv_G^2 + \frac{1}{2}I_G\omega^2$

For a hoop: $I_G = mR^2$

The velocity of center $G$: $v_G = \dot{s}$ where $s$ is the distance along the incline.

The angular velocity: $\omega = \dot{\theta}$

So: $E_c = \frac{1}{2}m\dot{s}^2 + \frac{1}{2}mR^2\dot{\theta}^2$

Potential energy (taking reference at the bottom of the incline):
$$E_p = mgh = mgs\sin\alpha$$

Mechanical energy:
$$E_m = E_c + E_p = \frac{1}{2}m\dot{s}^2 + \frac{1}{2}mR^2\dot{\theta}^2 + mgs\sin\alpha$$

\textbf{2. Differential equation using energy theorem:}

The theorem of variation of mechanical energy states:
$$\frac{dE_m}{dt} = P_{\text{non-conservative}}$$

The non-conservative power comes from friction at the contact point.

The power of friction force $\vec{T}$ is:
$$P = \vec{T} \cdot \vec{v}_g$$

Where $\vec{v}_g$ is the sliding velocity: $\vec{v}_g = \vec{V}_G + \vec{\omega} \times \vec{GI} = \dot{s}\vec{i} + \dot{\theta}\vec{k} \times (-R\vec{j}) = (\dot{s} + R\dot{\theta})\vec{i}$

So: $P = T(\dot{s} + R\dot{\theta})$

Now differentiate the mechanical energy:
$$\frac{dE_m}{dt} = m\dot{s}\ddot{s} + mR^2\dot{\theta}\ddot{\theta} + mg\dot{s}\sin\alpha$$

Equate to the non-conservative power:
$$m\dot{s}\ddot{s} + mR^2\dot{\theta}\ddot{\theta} + mg\dot{s}\sin\alpha = T(\dot{s} + R\dot{\theta})$$

We also have the equations from momentum theorems:
$$m\ddot{s} = mg\sin\alpha - T$$
$$I_G\ddot{\theta} = -RT \Rightarrow mR^2\ddot{\theta} = -RT$$

Substitute these into the energy equation:
$$\dot{s}(mg\sin\alpha - T) + R\dot{\theta}(-T) + mg\dot{s}\sin\alpha = T(\dot{s} + R\dot{\theta})$$

Simplify:
$$mg\dot{s}\sin\alpha - T\dot{s} - TR\dot{\theta} + mg\dot{s}\sin\alpha = T\dot{s} + TR\dot{\theta}$$
$$2mg\dot{s}\sin\alpha - T(\dot{s} + R\dot{\theta}) = T(\dot{s} + R\dot{\theta})$$
$$2mg\dot{s}\sin\alpha = 2T(\dot{s} + R\dot{\theta})$$
$$T = \frac{mg\dot{s}\sin\alpha}{\dot{s} + R\dot{\theta}}$$

This gives the friction force in terms of the motion. The differential equations are:
$$\ddot{s} = g\sin\alpha - \frac{g\dot{s}\sin\alpha}{\dot{s} + R\dot{\theta}}$$
$$\ddot{\theta} = -\frac{g\dot{s}\sin\alpha}{R(\dot{s} + R\dot{\theta})}$$
\end{correctionbox}

% ----------- EXERCISE 4 -------------
\section{}

We consider again the system from exercise II of series 5. The revolute joints at points $O$ and $C$ are assumed to be perfect. Initially $\theta = \theta_0$ and $\dot{\theta} = 0$.

\begin{enumerate}
    \item Determine the kinetic energy of the system with respect to $T_1$.
    \item Determine the potential energy of the system with respect to $T_1$. Deduce the mechanical energy of the system.
    \item By applying the angular momentum theorem to disk $S_2$ at point $C$, show that the rotation velocity of $S_2$ is constant.
    \item Using the theorem of variation of mechanical energy of the system, find the differential equation of motion, then deduce the period of small oscillations.
    \item Determine the period of small oscillations $T'$ in the case where the two solids are rigidly connected. Compare the two periods $T$ and $T'$.
\end{enumerate}

\begin{correctionbox}
Referring to Exercise 2 from Series 5, we have a system with rod $S_1$ and disk $S_2$.

\textbf{1. Kinetic energy:}

From Series 5, we found:
- Momentum: $\vec{P} = l(m_1/2 + m_2)\dot{\theta}\vec{j}$
- Angular momentum at $O_1$: $\vec{\sigma}_{O_1} = \left[l^2\left(\frac{m_1}{3} + m_2\right)\dot{\theta} + \frac{1}{2}m_2R^2\dot{\varphi}\right]\vec{k}_1$

The kinetic energy can be calculated using Koenig's theorem or directly from the definition.

For the rod $S_1$:
$$E_c(S_1) = \frac{1}{2}I_{O_1}\dot{\theta}^2 = \frac{1}{2} \cdot \frac{1}{3}m_1l^2\dot{\theta}^2 = \frac{1}{6}m_1l^2\dot{\theta}^2$$

For the disk $S_2$:
$$E_c(S_2) = \frac{1}{2}m_2v_C^2 + \frac{1}{2}I_C\dot{\varphi}^2$$

We have $v_C = l\dot{\theta}$ and $I_C = \frac{1}{2}m_2R^2$

So: $E_c(S_2) = \frac{1}{2}m_2l^2\dot{\theta}^2 + \frac{1}{4}m_2R^2\dot{\varphi}^2$

Total kinetic energy:
$$E_c = \frac{1}{6}m_1l^2\dot{\theta}^2 + \frac{1}{2}m_2l^2\dot{\theta}^2 + \frac{1}{4}m_2R^2\dot{\varphi}^2$$
$$= \left(\frac{m_1}{6} + \frac{m_2}{2}\right)l^2\dot{\theta}^2 + \frac{1}{4}m_2R^2\dot{\varphi}^2$$

\textbf{2. Potential energy:}

The potential energy is gravitational. Taking the pivot $O_1$ as reference:

For rod $S_1$: $E_p(S_1) = m_1g \cdot \frac{l}{2}\sin\theta = \frac{1}{2}m_1gl\sin\theta$

For disk $S_2$: $E_p(S_2) = m_2g \cdot l\sin\theta$

Total potential energy:
$$E_p = \left(\frac{m_1}{2} + m_2\right)gl\sin\theta$$

Mechanical energy:
$$E_m = \left(\frac{m_1}{6} + \frac{m_2}{2}\right)l^2\dot{\theta}^2 + \frac{1}{4}m_2R^2\dot{\varphi}^2 + \left(\frac{m_1}{2} + m_2\right)gl\sin\theta$$

\textbf{3. Constant rotation of disk $S_2$:}

Apply angular momentum theorem to disk $S_2$ about point $C$:

The external forces on $S_2$ are:
- Weight: $m_2\vec{g}$ acting at $C$ (zero moment about $C$)
- Reaction at the pivot $C$

If the joint at $C$ is perfect, the moment of the reaction about $C$ is zero.

Therefore: $\frac{d\vec{\sigma}_C(S_2)}{dt} = \vec{0}$

So $\vec{\sigma}_C(S_2)$ is constant.

But $\vec{\sigma}_C(S_2) = I_C\dot{\varphi}\vec{k} = \frac{1}{2}m_2R^2\dot{\varphi}\vec{k}$

Therefore $\dot{\varphi}$ is constant.

\textbf{4. Differential equation and period:}

Since $\dot{\varphi}$ is constant, the mechanical energy becomes:
$$E_m = \left(\frac{m_1}{6} + \frac{m_2}{2}\right)l^2\dot{\theta}^2 + \left(\frac{m_1}{2} + m_2\right)gl\sin\theta + \text{constant}$$

Differentiate with respect to time:
$$2\left(\frac{m_1}{6} + \frac{m_2}{2}\right)l^2\dot{\theta}\ddot{\theta} + \left(\frac{m_1}{2} + m_2\right)gl\dot{\theta}\cos\theta = 0$$

For small oscillations, $\sin\theta \approx \theta$, $\cos\theta \approx 1$:
$$\left(\frac{m_1}{3} + m_2\right)l^2\ddot{\theta} + \left(\frac{m_1}{2} + m_2\right)gl\theta = 0$$

This is a harmonic oscillator with frequency:
$$\omega^2 = \frac{\left(\frac{m_1}{2} + m_2\right)g}{\left(\frac{m_1}{3} + m_2\right)l}$$

Period: $T = 2\pi\sqrt{\frac{\left(\frac{m_1}{3} + m_2\right)l}{\left(\frac{m_1}{2} + m_2\right)g}}$

\textbf{5. Period for rigid connection:}

If the two solids are rigidly connected, $\dot{\varphi} = \dot{\theta}$, and the kinetic energy becomes:
$$E_c = \left(\frac{m_1}{6} + \frac{m_2}{2}\right)l^2\dot{\theta}^2 + \frac{1}{4}m_2R^2\dot{\theta}^2$$

The equation of motion:
$$\left[\left(\frac{m_1}{3} + m_2\right)l^2 + \frac{1}{2}m_2R^2\right]\ddot{\theta} + \left(\frac{m_1}{2} + m_2\right)gl\theta = 0$$

Frequency: $\omega'^2 = \frac{\left(\frac{m_1}{2} + m_2\right)gl}{\left(\frac{m_1}{3} + m_2\right)l^2 + \frac{1}{2}m_2R^2}$

Period: $T' = 2\pi\sqrt{\frac{\left(\frac{m_1}{3} + m_2\right)l^2 + \frac{1}{2}m_2R^2}{\left(\frac{m_1}{2} + m_2\right)gl}}$

Since the denominator in $T'$ is larger, $T' > T$. The rigid connection makes the system oscillate slower.
\end{correctionbox}

% ----------- EXERCISE 5 -------------
\section{}

Consider the system from exercise V of series 5 (hoop and a material point).

\begin{enumerate}
    \item Find the potential energy of the system.
    \item Calculate the power of external actions in two different ways.
    \item By applying the kinetic energy theorem to the system, find the differential equation of motion.
\end{enumerate}

\begin{correctionbox}
Referring to Exercise 5 from Series 5 (hoop with point mass).

\textbf{1. Potential energy:}

The system consists of:
- Hoop $D$ of mass $M$
- Point mass $A$ of mass $m$

Taking the horizontal through $O$ as reference:

For the hoop: $E_p(D) = MgR$ (constant)

For the point mass $A$: Its height is $y_A = R - R\cos\theta = R(1 - \cos\theta)$

So: $E_p(A) = mgR(1 - \cos\theta)$

Total potential energy:
$$E_p = MgR + mgR(1 - \cos\theta) = (M + m)gR - mgR\cos\theta$$

We can ignore the constant term for equations of motion:
$$E_p = -mgR\cos\theta$$

\textbf{2. Power of external actions:}

\textbf{First way:} From the definition $P = \sum \vec{F} \cdot \vec{v}$

External forces:
- Weight on hoop: $M\vec{g}$, acting at $C$, velocity $\vec{V}_C$
- Weight on point mass: $m\vec{g}$, acting at $A$, velocity $\vec{V}_A$
- Reaction at contact point $I$: $\vec{R}$, but $\vec{V}_I = \vec{0}$ (rolling without slipping)

So: $P = M\vec{g} \cdot \vec{V}_C + m\vec{g} \cdot \vec{V}_A$

We have $\vec{V}_C = R\dot{\phi}\vec{i}$, $\vec{V}_A = R(\dot{\phi} + \dot{\theta}\cos\theta)\vec{i} + R\dot{\theta}\sin\theta\vec{j}$

And $\vec{g} = -g\vec{j}$

So: $P = Mg\vec{j} \cdot (R\dot{\phi}\vec{i}) + mg\vec{j} \cdot [R(\dot{\phi} + \dot{\theta}\cos\theta)\vec{i} + R\dot{\theta}\sin\theta\vec{j}]$
$$= 0 + mg(-R\dot{\theta}\sin\theta) = -mgR\dot{\theta}\sin\theta$$

\textbf{Second way:} $P = -\frac{dE_p}{dt}$

From part 1: $E_p = -mgR\cos\theta$

So: $P = -\frac{d}{dt}(-mgR\cos\theta) = -mgR\sin\theta\dot{\theta}$

Both methods give the same result.

\textbf{3. Differential equation using kinetic energy theorem:}

The kinetic energy theorem states: $\frac{dE_c}{dt} = P_{\text{external}} + P_{\text{internal}}$

Since we have a rigid body system, $P_{\text{internal}} = 0$

From Series 5, we had the kinetic energy.

Let's derive it directly:

For the hoop: $E_c(D) = \frac{1}{2}Mv_C^2 + \frac{1}{2}I_C\omega^2 = \frac{1}{2}M(R\dot{\phi})^2 + \frac{1}{2}(MR^2)\dot{\phi}^2 = \frac{1}{2}MR^2\dot{\phi}^2 + \frac{1}{2}MR^2\dot{\phi}^2 = MR^2\dot{\phi}^2$

For the point mass: $E_c(A) = \frac{1}{2}mv_A^2$

We have $v_A^2 = [R(\dot{\phi} + \dot{\theta}\cos\theta)]^2 + [R\dot{\theta}\sin\theta]^2 = R^2[(\dot{\phi} + \dot{\theta}\cos\theta)^2 + (\dot{\theta}\sin\theta)^2]$
$$= R^2[\dot{\phi}^2 + 2\dot{\phi}\dot{\theta}\cos\theta + \dot{\theta}^2\cos^2\theta + \dot{\theta}^2\sin^2\theta] = R^2[\dot{\phi}^2 + 2\dot{\phi}\dot{\theta}\cos\theta + \dot{\theta}^2]$$

So: $E_c(A) = \frac{1}{2}mR^2(\dot{\phi}^2 + 2\dot{\phi}\dot{\theta}\cos\theta + \dot{\theta}^2)$

Total kinetic energy:
$$E_c = MR^2\dot{\phi}^2 + \frac{1}{2}mR^2(\dot{\phi}^2 + 2\dot{\phi}\dot{\theta}\cos\theta + \dot{\theta}^2)$$
$$= R^2\left[(M + \frac{m}{2})\dot{\phi}^2 + m\dot{\phi}\dot{\theta}\cos\theta + \frac{m}{2}\dot{\theta}^2\right]$$

Now, from the rolling without slipping condition, we have a relationship between $\phi$ and $\theta$.

Differentiate the kinetic energy:
$$\frac{dE_c}{dt} = R^2\left[2(M + \frac{m}{2})\dot{\phi}\ddot{\phi} + m\ddot{\phi}\dot{\theta}\cos\theta + m\dot{\phi}\ddot{\theta}\cos\theta - m\dot{\phi}\dot{\theta}^2\sin\theta + m\dot{\theta}\ddot{\theta}\right]$$

This equals the power: $-mgR\dot{\theta}\sin\theta$

This gives the differential equation of motion, which can be simplified using the constraints.

The final differential equation in $\theta$ is:
$$(M + m)R^2\ddot{\theta} + mgR\sin\theta = 0$$

For small oscillations: $\ddot{\theta} + \frac{mg}{(M + m)R}\theta = 0$

This is the equation of a harmonic oscillator.
\end{correctionbox}

\end{document}